% Lecture 2 slides. Ported from tmp/лекции/Lecture 2. Арифметика больших чисел.src.pptx
% (scripts/pptx2beamer.py + manual pass). Slide k of this deck = slide k of the PPTX.
% Changes against the PPTX:
%  * slides 3 and 9 were handwritten (pen) pages; they are typed here;
%  * slide 6: squaring c_t = 2 sum_{i=0}^{floor(t/2)} ... (the PPTX had ceil, wrong for odd t);
%  * slide 8: (1+2345)(67+8901) = 21038928 (the PPTX had 2818872);
%  * slide 13: 7 * 10000 (the PPTX had 7 * 1000).
\documentclass{ca-slides}

\lecture{2}{Арифметика больших чисел}
\author{}

\newcommand{\Len}{\mathcal{L}}

\begin{document}
\titleframe

% --- slide 2 ---
\begin{canvasframe}{Базовые операции арифметики}
  \begin{itemize}
    \item Сложение
    \item Умножение
    \begin{itemize}
      \item школьный метод
      \item быстрые алгоритмы: методы Карацубы, Тоома"--~Кука, Шёнхаге"--~Штрассена
    \end{itemize}
    \item Возведение в степень
    \begin{itemize}
      \item бинарный алгоритм возведения в степень
      \item оконный алгоритм возведения в степень
    \end{itemize}
    \item Целочисленное деление
    \begin{itemize}
      \item школьный метод
      \item алгоритм D "--- улучшенный школьный метод
    \end{itemize}
  \end{itemize}
\end{canvasframe}

% --- slide 3 (handwritten in the PPTX) ---
\begin{canvasframe}[shrink=18]{Оценки сложности}
  \small
  Пусть $g(x) > 0$ при всех достаточно больших $x \in G \subseteq \RR$.
  \begin{definition}[$\bigO$]
    $f(x) = \bigO(g(x))$ при $x \to \infty$, если $\exists\, C > 0, x_C > 0$:
    $\abs{f(x)} \le C g(x)$ $\forall x > x_C$;\\
    $f(x) = \bigO(g(x))$ при $x \to a$, если $\exists\, C > 0, \delta_C > 0$:
    $\abs{f(x)} \le C g(x)$ $\forall\, 0 < \abs{x - a} < \delta_C$,
    т.\,е. $\limsup_{x \to a} \abs{f(x)}/g(x) < \infty$.
  \end{definition}
  \begin{definition}[$o$]
    $f(x) = o(g(x))$ при $x \to \infty$, если $\forall \varepsilon > 0\ \exists N_\varepsilon$:
    $\abs{f(x)} \le \varepsilon g(x)$ $\forall x > N_\varepsilon$, т.\,е. $\lim f/g = 0$.
  \end{definition}
  \begin{definition}[$\Omega$, $\Theta$, $\sim$; по Кнуту]
    $f = \Omega(g) \Leftrightarrow g = \bigO(f)$;\quad
    $f = \Theta(g) \Leftrightarrow f = \bigO(g)$ и $g = \bigO(f)$;\quad
    $f \sim g \Leftrightarrow \lim_{x \to a} f(x)/g(x) = 1$.
  \end{definition}
\end{canvasframe}

% --- slide 4 ---
\begin{canvasframe}[shrink=16]{Представление больших чисел}
  \begin{itemize}
    \item Основание системы счисления будем обозначать $b$ или $B$
    \item \emph{Представлением} неотрицательного \textbf{целого} числа $x$ называется
      кратчайшая последовательность $\{x_i\}$, $0 \le x_i < B$:\quad
      $x = \sum_{i=0}^{D-1} x_i B^i$
    \item \emph{Сбалансированным} представлением называется кратчайшая последовательность
      $\{\hat x_i\}$, $-\floor{B/2} \le \hat x_i < \floor{B/2}$:\quad
      $x = \sum_{i=0}^{D-1} \hat x_i B^i$
    \item Количество цифр $x_i$ в представлении числа по основанию $B$ называется
      $B$-\emph{длиной} и равно $\Len(x) = \floor{\log_B x} + 1$
  \end{itemize}
  \aside{единственно ли сбалансированное представление?}
  \aside{как перевести из одного представления в другое?}
\end{canvasframe}

% --- slide 5 ---
\begin{canvasframe}{Сложение больших чисел}
  \begin{itemize}
    \item Пусть $a, b \in \NN$, тогда
      \[\max(\Len(a), \Len(b)) \le \Len(a + b) \le \max(\Len(a), \Len(b)) + 1\]
  \end{itemize}
  \aside{сформулируйте алгоритм сложения двух чисел, представленных в сбалансированной форме}
  \penspace
\end{canvasframe}

% --- slide 6 ---
\begin{canvasframe}[shrink=18]{Умножение больших чисел}
  \small
  \begin{itemize}
    \item Пусть $a, b \in \NN$, $a = \sum_{i=0}^{M-1} a_i B^i$, $b = \sum_{i=0}^{N-1} b_i B^i$,
      $c = ab$, тогда $c_t = \sum_{i+j=t} a_i b_j$
    \item Длина произведения: $\Len(a) + \Len(b) - 1 \le \Len(c) \le \Len(a) + \Len(b)$
    \item Для квадрата $c = a \cdot a$: $c_t = \sum_{i=0}^{t} a_i a_{t-i}$; ввиду зеркальной
      симметрии
      \[c_t = 2 \sum_{i=0}^{\floor{t/2}} a_i a_{t-i} - \mu_t,\qquad
        \mu_t = \begin{cases} 0, & t \text{ нечётно},\\ a_{t/2}^2, & t \text{ чётно}\end{cases}\]
    \item Следовательно, возведение в квадрат почти \textbf{в два раза быстрее} умножения
    \item Сложность школьного алгоритма умножения: $\bigO(N^2)$ умножений цифр
  \end{itemize}
  \aside{откуда эта формула?}
  \aside{умножение не может быть в два раза «сложнее» возведения в квадрат "--- почему?}
\end{canvasframe}

% --- slide 7 ---
\begin{canvasframe}[shrink=16]{Умножение: алгоритм Карацубы}
  \small
  \begin{itemize}
    \item Пусть $a, b \in \NN$, $a = \sum_{i=0}^{N-1} a_i B^i$, $b = \sum_{i=0}^{N-1} b_i B^i$
    \item Не теряя общности, $N = 2k$, $K = B^k$, т.\,е. $a, b < K^2$
    \item Пусть $a = a_1 K + a_0$, $b = b_1 K + b_0$ (числа $a$ и $b$ записаны по основанию $K$)
    \item Тогда $ab = a_1 b_1 K^2 + (a_1 b_0 + a_0 b_1) K + a_0 b_0$
    \item Учитывая $a_1 b_0 + a_0 b_1 = (a_1 + a_0)(b_1 + b_0) - a_1 b_1 - a_0 b_0$,
      \[ab = a_1 b_1 K^2 + \bigl((a_1 + a_0)(b_1 + b_0) - a_1 b_1 - a_0 b_0\bigr) K + a_0 b_0\]
    \item Умножение $N$-значных чисел сводится к \textbf{трём} умножениям чисел размера $N/2$,
      сложениям и сдвигам
    \item Сложность умножения: $\bigO(N^{\log_2 3})$, $\log_2 3 \approx 1{,}585 < 2$
  \end{itemize}
  \aside{как получено выражение сложности рекуррентной схемы?}
\end{canvasframe}

% --- slide 8 ---
\begin{canvasframe}[shrink=12]{Умножение: алгоритм Карацубы}
  \small
  \begin{itemize}
    \item Вычислим произведение двух целых размера $\bigO(10^8)$, $K = 10^4$:
      \[a = 12345 = 0001 \cdot 10^4 + 2345,\qquad b = 678901 = 0067 \cdot 10^4 + 8901\]
    \item Три умножения целых размера $\bigO(10^4)$:
      \begin{align*}
        r_2 &= 1 \cdot 67 = 67, \qquad r_0 = 2345 \cdot 8901 = 20872845,\\
        r_1 &= (1 + 2345)(67 + 8901) - r_2 - r_0 = 21038928 - 67 - 20872845 = 166016
      \end{align*}
    \item Собираем результат:
      $ab = r_2 \cdot 10^8 + r_1 \cdot 10^4 + r_0 = 67 \cdot 10^8 + 166016 \cdot 10^4 + 20872845
      = 8381032845$
    \item Дабы «увеличить выигрыш», повторяем трюк для двух умножений чисел размера
      $\bigO(10^4)$: $2345 \cdot 8901$ и $(1 + 2345)(67 + 8901)$ "--- «меняем» каждое на три
      умножения чисел размера $\bigO(10^2)$
  \end{itemize}
\end{canvasframe}

% --- slide 9 (handwritten in the PPTX) ---
\begin{canvasframe}[shrink=14]{Основная теорема о рекуррентных соотношениях}
  \footnotesize
  \begin{theorem}
    Пусть $T(n) = a\,T(n/b) + f(n)$, $c_{\mathrm{crit}} = \log_b a$ "--- критическая экспонента.
    \begin{enumerate}
      \item Если $f(n) = \bigO(n^c)$, $c < c_{\mathrm{crit}}$, то $T(n) = \Theta(n^{c_{\mathrm{crit}}})$.
      \item Если $f(n) = \Theta(n^{c} \log^k n)$, $c = c_{\mathrm{crit}}$, $k \ge 0$, то
        $T(n) = \Theta(n^{c} \log^{k+1} n)$.
      \item Если $f(n) = \Omega(n^c)$, $c > c_{\mathrm{crit}}$, и $a f(n/b) \le k f(n)$ для
        некоторого $k < 1$ (условие регулярности), то $T(n) = \Theta(f(n))$.
    \end{enumerate}
  \end{theorem}
  \begin{columns}[T]
    \column{0.33\textwidth}
      \textcolor{caRed}{1. разбиение «проще» подзадач}\\
      $T = 10\,T(n/3) + 42 n^2$:\\ $c = 2 = \log_3 9 < \log_3 10$\\
      $\Rightarrow T = \Theta(n^{\log_3 10})$
    \column{0.33\textwidth}
      \textcolor{caRed}{2. разбиение $\sim$ подзадачам}\\
      $T = 2\,T(n/4) + 23\sqrt n$:\\ $c_{\mathrm{crit}} = \log_4 2 = \frac12 = c$, $k = 0$\\
      $\Rightarrow T = \Theta(\sqrt n \log n)$
    \column{0.33\textwidth}
      \textcolor{caRed}{3. разбиение «сложнее» подзадач}\\
      $T = 4\,T(n/3) + n^2$: $c = 2 > \log_3 4$;\\
      $4 (n/3)^2 = \frac49 n^2$, $k = \frac49 < 1$\\
      $\Rightarrow T = \Theta(n^2)$
  \end{columns}
\end{canvasframe}

% --- slide 10 ---
\begin{canvasframe}[shrink=16]{Умножение: алгоритм Тоома"--~Кука}
  \small
  \begin{itemize}
    \item В алгоритме Карацубы задача размера $N$ делилась на три задачи размера $N/2$
    \item \alert{Вопрос}: как получить асимптотически более эффективный алгоритм?
    \item Для фиксированного $D \in \NN$ пусть $N = Dk$, $K = B^k$, т.\,е. $a, b < K^D$
      (числа $a$ и $b$ делим на $D$ частей)
    \item Сомножители представим в виде полиномов
      \[a(x) = a_0 + a_1 x + \cdots + a_{D-1} x^{D-1},\qquad
        b(x) = b_0 + b_1 x + \cdots + b_{D-1} x^{D-1},\]
      где коэффициенты выбирают так, чтобы $a(K) = a$, $b(K) = b$
    \item $a = 126386_{10} = 1034321_7$:\quad
      $a_{10}(x) = 6 + 8x + 3x^2 + 6x^3 + 2x^4 + x^5$,\quad
      $a_7(x) = 1 + 2x + 3x^2 + 4x^3 + 3x^4 + 0x^5 + x^6$,\quad
      $a_{100}(x) = 86 + 63x + 12x^2$
  \end{itemize}
  \aside{пусть задан многочлен степени $D-1$ с неизвестными коэффициентами: сколько его значений
    в различных точках нужно, чтобы определить коэффициенты?}
\end{canvasframe}

% --- slide 11 ---
\begin{canvasframe}[shrink=8]{Умножение: алгоритм Тоома"--~Кука}
  \small
  \begin{itemize}
    \item $f(x) = a_0 + a_1 x + a_2 x^2 + \cdots + a_{D-1} x^{D-1}$; оракул выдаёт значения в любой точке:
      \[\left\{\begin{aligned}
        a_0 + 0a_1 + 0a_2 + \cdots + 0a_{D-1} &= f(0)\\
        a_0 + a_1 + a_2 + \cdots + a_{D-1} &= f(1)\\
        a_0 - a_1 + a_2 - \cdots + (-1)^{D-1} a_{D-1} &= f(-1)\\
        \ldots&
      \end{aligned}\right.\]
    \item Система уравнений с $D$ неизвестными имеет единственное решение, если различных точек
      не меньше $D$
    \item $f(x) = a_0 + a_1 x + a_2 x^2 + a_3 x^3 + a_4 x^4$, $\deg f = 4$:\quad
      $f(-2) = 9$, $f(-1) = 3$, $f(0) = 5$, $f(1) = 15$, $f(2) = 57$.\quad $f = {?}$
  \end{itemize}
  \aside{систему можно переписать в матричном виде; определитель матрицы называется
    определителем Вандермонда}
\end{canvasframe}

% --- slide 12 ---
\begin{canvasframe}{Умножение: алгоритм Тоома"--~Кука}
  \[\left\{\begin{aligned}
    a_0 &= f(0) = 5\\
    a_0 + a_1 + a_2 + a_3 + a_4 &= f(1) = 15\\
    a_0 - a_1 + a_2 - a_3 + a_4 &= f(-1) = 3\\
    a_0 + 2a_1 + 4a_2 + 8a_3 + 16a_4 &= f(2) = 57\\
    a_0 - 2a_1 + 4a_2 - 8a_3 + 16a_4 &= f(-2) = 9
  \end{aligned}\right.\]
  \[f(x) = 5 + 4x + 3x^2 + 2x^3 + x^4\]
  \penspace[0.15\textheight]
\end{canvasframe}

% --- slide 13 ---
\begin{canvasframe}[shrink=16]{Умножение: алгоритм Тоома"--~Кука}
  \begin{itemize}
    \item Представим сомножители в виде полиномов $a(x)$, $b(x)$ степени $D-1$, $a(K) = a$, $b(K) = b$
    \item Пусть $c(x) = a(x) b(x)$ и $c = ab$. \textbf{Верно ли} $c(K) = a(K) b(K) = ab = c$?
    \item \textbf{Какова степень} полинома $c(x)$?
    \item $K = 10$, $a = 157$, $b = 768$; $a(x) = 7 + 5x + x^2$, $b(x) = 8 + 6x + 7x^2$
    \item $c(x) = 56 + 82x + 87x^2 + 41x^3 + 7x^4$, \quad $c = ab = 157 \cdot 768 = 120576$
    \item $c(10) = 56 + 82 \cdot 10 + 87 \cdot 100 + 41 \cdot 1000 + 7 \cdot 10000 = 120576$ "--- успех!
  \end{itemize}
  \aside{а это как доказать?}
\end{canvasframe}

% --- slide 14 ---
\begin{canvasframe}[shrink=28]{Умножение: алгоритм Тоома"--~Кука}
  \small
  Метод основан на том, что \textbf{произведение многочленов} $a(x)$ и $b(x)$ степени не выше
  $D - 1$, $c(x) = a(x) b(x)$, \textbf{полностью определяется значениями в $2D - 1$ точках}.
  \begin{enumerate}
    \item {[Подготовка]} Сформировать два \textbf{символических} многочлена $a(x)$, $b(x)$ степени $D-1$.
    \item {[Вычисление]} Для каждого $t \in [1-D, D-1]$ выразить $c(t) = a(t) b(t)$
      \textbf{символически}, т.\,е. через коэффициенты $a$ и $b$.
    \item {[Восстановление]} Составить и решить \textbf{символически} систему из $2D-1$ уравнений
      относительно $c_k$: $c(t) = \sum_{k=0}^{2D-2} c_k t^k$, $t \in [1-D, D-1]$.
      \textbf{N.B.} Величины $c(t)$ "--- символические выражения («формулы») этапа 2.
    \item {[Формирование результата]} Выписать выражения каждого $c_k$ через коэффициенты $a$ и $b$.
  \end{enumerate}
  Сложность умножения: $\bigO(N^{\log_D (2D-1)})$
  \aside{при $D = 3$: $\log_3 5 \approx 1{,}465 < \log_2 3 < 2$ "--- асимптотически лучше Карацубы}
  \aside{на практике $c(t)$ сразу вычисляют, выполняя $2D-1$ умножений чисел размера $N/D$}
\end{canvasframe}

% --- slide 15 ---
\begin{canvasframe}[shrink=20]{Умножение: алгоритм Тоома"--~Кука}
  \small
  Пример $D = 3$: $a(x) = a_0 + a_1 x + a_2 x^2$, $b(x) = b_0 + b_1 x + b_2 x^2$.
  \begin{itemize}
    \item {[Подготовка]}
      \[\begin{aligned}
        a(-2) &= a_0 - 2a_1 + 4a_2, & b(-2) &= b_0 - 2b_1 + 4b_2,\\
        a(-1) &= a_0 - a_1 + a_2, & b(-1) &= b_0 - b_1 + b_2,\\
        a(0) &= a_0, & b(0) &= b_0,\\
        a(1) &= a_0 + a_1 + a_2, & b(1) &= b_0 + b_1 + b_2,\\
        a(2) &= a_0 + 2a_1 + 4a_2, & b(2) &= b_0 + 2b_1 + 4b_2
      \end{aligned}\]
    \item {[Вычисление]} $t_0 = c(-2) = a(-2) b(-2)$, $t_1 = c(-1)$, $t_2 = c(0) = a_0 b_0$,
      $t_3 = c(1)$, $t_4 = c(2)$.
      Все «длинные» умножения выполняются на этом этапе: \textbf{пять умножений} для $D = 3$
  \end{itemize}
\end{canvasframe}

% --- slide 16 ---
\begin{canvasframe}[shrink=38]{Умножение: алгоритм Тоома"--~Кука}
  \small
  \begin{itemize}
    \item {[Восстановление]} Система для $D = 3$:
      \[\left\{\begin{aligned}
        c_0 &= c(0) = t_2\\
        c_0 + c_1 + c_2 + c_3 + c_4 &= c(1) = t_3\\
        c_0 - c_1 + c_2 - c_3 + c_4 &= c(-1) = t_1\\
        c_0 + 2c_1 + 4c_2 + 8c_3 + 16c_4 &= c(2) = t_4\\
        c_0 - 2c_1 + 4c_2 - 8c_3 + 16c_4 &= c(-2) = t_0
      \end{aligned}\right.\]
      Решая её \textbf{символически} (не подставляя значения $t_i$!), получаем
      \[\begin{aligned}
        c_0 &= t_2, & c_1 &= \tfrac{t_0}{12} - \tfrac{2t_1}{3} + \tfrac{2t_3}{3} - \tfrac{t_4}{12},\\
        c_2 &= -\tfrac{t_0}{24} + \tfrac{2t_1}{3} - \tfrac{5t_2}{4} + \tfrac{2t_3}{3} - \tfrac{t_4}{24}, &
        c_3 &= -\tfrac{t_0}{12} + \tfrac{t_1}{6} - \tfrac{t_3}{6} + \tfrac{t_4}{12},\\
        c_4 &= \tfrac{t_0}{24} - \tfrac{t_1}{6} + \tfrac{t_2}{4} - \tfrac{t_3}{6} + \tfrac{t_4}{24}
      \end{aligned}\]
      Эти операции «дешёвые» по сравнению с «длинными» умножениями этапа 2
    \item {[Формирование результата]} Подставляя цифры, вычисляем $c_k$; с точностью до
      переносов они совпадают с цифрами произведения $c = ab$
  \end{itemize}
\end{canvasframe}

% --- slide 17 ---
\begin{canvasframe}[shrink=6]{Умножение: алгоритм Тоома"--~Кука}
  \small
  $a = 549238$, $b = 807541$, $K = 100$:\quad $a(x) = 38 + 92x + 54x^2$, $b(x) = 41 + 75x + 80x^2$.
  \[\begin{aligned}
    t_0 &= (38 - 2 \cdot 92 + 4 \cdot 54)(41 - 2 \cdot 75 + 4 \cdot 80) = 70 \cdot 211 = 14770,\\
    t_1 &= (38 - 92 + 54)(41 - 75 + 80) = 0, \qquad t_2 = 38 \cdot 41 = 1558,\\
    t_3 &= (38 + 92 + 54)(41 + 75 + 80) = 36064,\\
    t_4 &= (38 + 2 \cdot 92 + 4 \cdot 54)(41 + 2 \cdot 75 + 4 \cdot 80) = 223818
  \end{aligned}\]
  $c_0 = 1558$, $c_1 = 6622$, $c_2 = 12154$, $c_3 = 11410$, $c_4 = 4320$:\quad
  $c(x) = 1558 + 6622x + 12154x^2 + 11410x^3 + 4320x^4$.

  Переносы от $c_0$ до $c_4$ (с дополнительным разрядом $c_5$) дают цифры по основанию $100$
  (младшие "--- первыми): $c = (58, 37, 20, 32, 35, 44)$, т.\,е. $ab = 443532203758$.
\end{canvasframe}

% --- slide 18 ---
\begin{canvasframe}{Умножение: алгоритм Тоома"--~Кука}
  \begin{itemize}
    \item Метод Тоома"--~Кука позволяет достичь сложности мультипликативной части
      $\bigO(N^{1+\varepsilon})$ $\forall \varepsilon > 0$
    \item Однако уменьшение количества умножений достигается за счёт увеличения количества сложений
    \item Более того, коэффициенты в символических формулах (этап [Восстановление]) растут, что
      усложняет восстановление коэффициентов $c_k$
    \item На практике эти недостатки для «реальных» размеров сомножителей оказываются существенными
    \item Таким образом, алгоритм представляет в первую очередь теоретический интерес
  \end{itemize}
\end{canvasframe}
\end{document}
